% !TeX root = ../report.tex
% Chapter 1: Introduction

\chapter{บทนำ}
% \thaijustify{
%   [ภาพรวมของบทนี้]
% }
\section{ที่มาและความสำคัญของปัญหา}
\thaijustify{
การออกแบบประสบการณ์ผู้ใช้และส่วนต่อประสานผู้ใช้ (UX/UI Design) โดยปกติมีลักษณะเป็นกระบวนการที่วนซ้ำและยึดผู้ใช้เป็นศูนย์กลาง ทว่าในทางปฏิบัตินักออกแบบมักถูกจำกัดด้วยเวลา ประสบการณ์ และทรัพยากร กระบวนการออกแบบเริ่มตั้งแต่การรับโจทย์ รวบรวมความต้องการ ทำความเข้าใจผู้ใช้ นิยามปัญหา ระดมความคิด สร้างต้นแบบ และทดสอบ ในขณะเดียวกันยังต้องทำงานร่วมกับผู้มีส่วนได้ส่วนเสียอื่น ในกระบวนการสร้าง wireframe หรือ prototype มักถูกสร้างขึ้นด้วยมือ กระบวนการนี้จึงใช้แรงงานและเวลามาก~\cite{feng2023designing} ทำให้นักออกแบบสำรวจทางเลือกได้ไม่กว้าง
และเก็บข้อเสนอแนะจากผู้ใช้ได้ไม่เพียงพอ~\cite{suh2025storyensemble}
โอกาสที่ข้อเสนอแนะจะถูกนำมาปรับปรุงงานก่อนส่งมอบจึงลดลง สิ่งนี้ก่อให้เกิดช่องว่างระหว่างผลิตภัณฑ์ที่ส่งมอบกับความต้องการที่แท้จริงของผู้ใช้ ซึ่งเป็นผลมาจากการตัดสินใจเชิงออกแบบที่ถูกกำหนดโดยแรงกดดันด้านเวลามากกว่าความต้องการของผู้ใช้จริง}

\thaijustify{
  งานวิจัยที่ผ่านมาชี้ว่า Generative AI สามารถช่วยลดภาระทางด้านเวลาและแรงงานในการสร้าง prototype และลดปัญหาคอขวดของการทดสอบได้ อย่างไรก็ตาม เครื่องมือที่มีอยู่ในปัจจุบันยังคงกระจัดกระจายและส่วนใหญ่มุ่งเน้นเฉพาะงานใดงานหนึ่งเป็นหลัก ยังไม่ถูกนำมาใช้ตลอดทั้ง UX/UI Development Life Cycle \cite{chung2025ureka, lu2022bridging}
}

% \thaijustify{
%   งานวิจัยที่ผ่านมาชี้ให้เห็นว่า Generative AI สามารถช่วยลดช่องว่างดังกล่าวได้ โดยทำให้แต่ละรอบของการวนซ้ำใช้เวลาน้อยลงและส่งผลให้มีการวนซ้ำอย่างสม่ำเสมอในกระบวนการพัฒนาผลิตภัญฑ์เป็นไปได้ในทางปฏิบัติ ผลการศึกษาแสดงให้เห็นว่า Generative AI สามารถสร้าง wireframe ได้โดยอัตโนมัติจากการอธิบายความต้องการทางการออกแบบด้วยภาษาธรรมชาติ \cite{feng2023designing} และยังสามารถสร้าง high-fidelity prototype ได้ทั้งระดับภาพรวมทั้งระบบและระดับ component \cite{yuan2024towards} ซึ่งช่วยลดต้นทุนด้านเวลาและแรงงานของการออกแบบในแต่ละรอบ นอกจากการทำให้การวนซ้ำทดสอบเกิดบ่อยขึ้นแล้ว สภาพแวดล้อมการออกแบบแบบ forward-backward propagation ยังเปิดโอกาสให้ทีมผู้พัฒนาย้อนกลับไปทบทวนและปรับแก้การตัดสินใจก่อนหน้าได้อย่างเป็นพลวัตตลอดหลายรอบการทำงาน แทนที่จะดำเนินไปข้างหน้าเพียงอย่างเดียว \cite{suh2025storyensemble} แนวทางนี้สอดคล้องกับงานวิจัยด้านการใช้งาน (usability) ที่เป็นที่ยอมรับ ซึ่งแสดงให้เห็นว่าการทดสอบบ่อยครั้งด้วยกลุ่มตัวอย่างขนาดเล็กให้ insight ที่นำไปปรับใช้ได้จริงมากกว่าการทดสอบขนาดใหญ่ที่ทำน้อยครั้งกว่า \cite{nielsen2000five} เมื่อพิจารณาโดยรวม งานวิจัยเหล่านี้ชี้ว่า Generative AI สามารถช่วยลดภาระทางด้านเวลาและแรงงานในการสร้าง prototype และปัญหาคอขวดของการทดสอบได้ อย่างไรก็ตาม เครื่องมือที่มีอยู่ในปัจจุบันยังคงกระจัดกระจายและส่วนใหญ่มุ่งเน้นเฉพาะงานใดงานหนึ่งเป็นหลัก ยังไม่ถูกนำมาใช้ตลอดทั้ง UX/UI DLC \cite{chung2025ureka, lu2022bridging}}

\thaijustify{
เพื่อเพิ่มประสิทธิภาพ UX/UI DLC โครงงานนี้ให้ความสำคัญกับกระบวนการ comparative usability testing เป็นพิเศษ ซึ่งการทำการทดสอบซ้ำหลายรอบ นักออกแบบจำเป็นต้องจัดเตรียม two-alternative UI, รวบรวมและจัดระเบียบข้อเสนอแนะจากผู้ใช้, ตลอดจนตีความผลลัพธ์ที่ได้ก่อนเริ่มการวนซ้ำรอบถัดไปด้วยตนเอง เพื่อสนับสนุนกระบวนการดังกล่าว ผู้พัฒนาจึงนำเสนอแพลตฟอร์มที่\mbox{บูรณาการ} Generative AI ที่มีอยู่เข้ากับกระบวนการ comparative usability testing กล่าวคือ นักออกแบบสามารถอธิบายความต้องการทางการออกแบบจากนั้นแพลตฟอร์มจะสร้าง two-alternative UI สำหรับใช้ในการทดสอบและผู้เข้าร่วมทดสอบสามารถให้ข้อเสนอแนะเชิงบริบทด้วยการเลือก component ที่ต้องการให้ข้อเสนอแนะ และข้อเสนอแนะเหล่านี้จะถูกนำไปรวบรวม กลั่นกรอง โดยปัญญาประดิษฐ์แล้วนักออกแบบตรวจสอบและยืนยันเพื่อสร้าง two-alternative UI ใหม่สำหรับการทำ comparative usability testing ในรอบถัดไป}

\thaijustify{
นอกจากนี้ แพลตฟอร์มดังกล่าวยังมุ่งสร้างขั้นตอนการทำงานร่วมกันระหว่างมนุษย์และปัญญาประดิษฐ์อย่างสมดุล (human-in-the-loop) โดยปัญญาประดิษฐ์จะสรุปและวิเคราะห์ข้อเสนอแนะที่ได้รับจากรอบการทดสอบ เพื่อระบุรูปแบบ ข้อมูลเชิงลึก และประเด็นสำคัญ จากนั้นนักออกแบบจะเข้ามามีบทบาทในการพิจารณาผลการวิเคราะห์ที่ปัญญาประดิษฐ์สร้างขึ้น ใช้การตัดสินใจอย่างมืออาชีพและมีวิจารณญาณในการกำหนดทิศทางการออกแบบ \cite{chung2025ureka} เมื่อนักออกแบบตรวจสอบและยืนยันแนวทางแล้ว Generative AI จะดำเนินการตามที่นักออกแบบกำหนดเพื่อสร้าง two-alternative UI สำหรับรอบการทดสอบถัดไป อีกทั้งยังสร้างรายงานการออกแบบ ที่บันทึกว่าได้รับข้อเสนอแนะใดบ้าง ข้อเสนอแนะนั้นได้รับการตีความโดยปัญญาประดิษฐ์และนักออกแบบตัดสินใจอย่างไร เพื่อเป็นเหตุผลประกอบที่ชัดเจนและมีข้อมูลสนับสนุนสำหรับการสื่อสารกับผู้มีส่วนได้ส่วนเสีย โดยสรุป โครงงานนี้มีเป้าหมายเพื่อแสดงให้เห็นว่าแพลตฟอร์มที่เสริมศักยภาพด้วยการใช้ปัญญาประดิษฐ์ร่วมกับมนุษย์ใน UX/UI DLC สามารถลดระยะเวลาและภาระงานที่ต้องทำด้วยมือ เพิ่มความถี่ของวงรอบการทดสอบ และยกระดับคุณภาพโดยรวมของผลิตภัณฑ์ที่ถูกพัฒนา ตลอดจนเป็นเครื่องมือที่ใช้ตรวจสอบย้อนกลับการตัดสินใจตลอดทั้งกระบวนการได้}

\section{วัตถุประสงค์ของโครงงาน}
\thaijustify{
  \begin{enumerate}

    \item เพื่อให้นักออกแบบสามารถอธิบายว่าข้อเสนอแนะจากผู้ใช้ส่งผลต่อการตัดสินใจเชิงออกแบบอย่างไร และสามารถนำเหตุผลประกอบที่สืบกลับได้ไปใช้สร้าง persona artifact สำหรับผู้มีส่วนได้ส่วนเสียแต่ละกลุ่ม เพื่อใช้เป็นเหตุผลรองรับโดยมีข้อมูลสนับสนุนผลลัพธ์การออกแบบในขั้นสุดท้าย

    \item เพื่อลดภาระงานที่ต้องทำด้วยมือและระยะเวลาที่นักออกแบบต้องใช้ในการสร้าง two-alternative UI prototype สำหรับ comparative usability testing

    \item เพื่อเสริมศักยภาพให้นักออกแบบสามารถดำเนินการทดสอบกับผู้ใช้ได้บ่อยครั้งและสม่ำเสมอภายใต้กรอบเวลาส่งมอบงานที่จำกัด โดยอาศัย Generative AI ในการปรับปรุงการออกแบบบนพื้นฐานของข้อเสนอแนะที่ผ่านการตรวจสอบยืนยันแล้ว และนำไปใช้สำหรับ comparative usability testing รอบถัดไป

    \item เพื่อให้ผู้เข้าร่วมทดสอบ comparative usability testing สามารถให้ข้อเสนอแนะเชิงบริบทมากขึ้นโดยใช้ความพยายามน้อยลง

    \item เพื่อลดภาระงานของนักออกแบบในการวิเคราะห์ข้อเสนอแนะจากผู้ใช้ที่มีปริมาณมาก
  \end{enumerate}
}
\section{ขอบเขตของโครงงาน}

\thaijustify{
  \begin{enumerate}
    \item ใช้โมเดล Generative AI ที่มีอยู่ในปัจจุบันเพื่อสร้าง two-alternative UI prototype จากการอธิบายความต้องการทางการออกแบบด้วยภาษาธรรมชาติ
    \item จัดเตรียมสภาพแวดล้อมสำหรับให้ผู้เข้าร่วมทดสอบ comparative usability testing ตรวจสอบ เปรียบเทียบ และให้ข้อเสนอแนะเชิงบริบท
    \item ใช้แบบโมเดล AI ที่มีอยู่ในปัจจุบันในการรวบรวมข้อเสนอแนะและระบุรูปแบบ ข้อมูลเชิงลึก และประเด็นสำคัญเพื่อพิจารณา โดยยังคงให้นักออกแบบเป็นผู้ตัดสินใจ (human-in-the-loop) ก่อนเริ่มการทำ comparative usability testing ในรอบถัดไป
    \item ใช้การวิเคราะห์ที่ได้จาก AI ในการเชื่อมโยงข้อเสนอแนะของผู้ใช้เข้ากับวิวัฒนาการของการออกแบบและสร้างรายงาน
  \end{enumerate}
}

\subsection{คุณสมบัติของซอฟต์แวร์ (Software Features)}
\thaijustify{
  \todo{รายการฟีเจอร์หลักของแพลตฟอร์ม แยกตามส่วนของระบบ (Authentication และสิทธิ์ตามบทบาท, Project setup และ Design intent, AI Design Generation, Testing space, Feedback analysis, Stakeholder report)}
}

\subsection{กลุ่มเป้าหมาย}
\thaijustify{
  \todo{ผู้ใช้แต่ละบทบาท (Designer, Usability Tester, Client, Developer, Project Manager) และลักษณะของแต่ละกลุ่ม}
}

\subsection{สภาพแวดล้อม}
\thaijustify{
  แพลตฟอร์มพัฒนาเป็นเว็บแอปพลิเคชันที่ใช้งานผ่านเว็บเบราว์เซอร์ โดยไม่ต้องติดตั้งโปรแกรมเพิ่มเติม ผู้ใช้ทุกบทบาทใช้งานผ่านเบราว์เซอร์บนคอมพิวเตอร์ ส่วนผู้เข้าร่วมทดสอบเข้าสู่เซสชันการทดสอบด้วยรหัสหรือลิงก์ สภาพแวดล้อมของระบบมีรายละเอียดดังนี้
  \begin{itemize}
    \item \textbf{ฝั่งผู้ใช้:} เว็บเบราว์เซอร์รุ่นปัจจุบัน ส่วนหน้าจอที่ AI สร้างขึ้นรองรับการแสดงผลสามแบบ ได้แก่ desktop, tablet และ mobile ตามที่กำหนดในแต่ละโปรเจกต์
    \item \textbf{ฝั่งเซิร์ฟเวอร์:} Node.js รุ่น 20.19 ขึ้นไป ใช้ฐานข้อมูล PostgreSQL 17 และ Redis 7 ซึ่งรันผ่าน Docker
    \item \textbf{บริการ AI:} เรียกใช้โมเดลภาษาขนาดใหญ่ผ่าน CPE AI Gateway ของภาควิชาวิศวกรรมคอมพิวเตอร์ ซึ่งเข้าถึงได้ผ่านเครือข่าย VPN ของมหาวิทยาลัยเท่านั้น
    \item \textbf{การนำระบบขึ้นใช้งาน:} \todo{ระบุเซิร์ฟเวอร์หรือบริการที่ใช้ deploy ระบบ (ใน repository ยังไม่มีการตั้งค่า deploy)}
  \end{itemize}
}

\section{แผนการดำเนินงาน}
\subsection{Task Breakdown}
\begin{enumerate}
  \item \textbf{Documentation}
    \begin{itemize}
      \item ศึกษาค้นคว้าหัวข้อโครงงาน
      \item จัดทำเอกสาร Project Idea
      \item เอกสารและงานวิจัยที่เกี่ยวข้อง
      \item จัดทำเอกสาร Proposal report
      \item นำเสนอ Proposal
      \item จัดทำเอกสาร Progress report ภาคเรียนที่ 1
      \item นำเสนอ Progress ภาคเรียนที่ 1
      \item จัดทำเอกสาร Progress report ภาคเรียนที่ 2
      \item นำเสนอ Progress ภาคเรียนที่ 2
      \item จัดทำเอกสาร Final Report
      \item นำเสนอ Final Defense
    \end{itemize}

  \item \textbf{Requirements Analysis \& System Design}
    \begin{itemize}
      \item สัมภาษณ์ผู้มีส่วนได้ส่วนเสีย
      \item ระบุ Requirements
      \item ระบุ Functional \& Non-Functional Requirements
      \item จัดทำ Use Cases \& Sequence Diagram
      \item วางโครงสร้าง System Architecture
      \item วางโครงสร้าง AI Workflow \& Prompt Architecture
      \item ออกแบบ Database
      \item ออกแบบ API Design
      \item ออกแบบ UI/UX Design (User Flow \& Wireframe)
      \item System prototype (PoC)
    \end{itemize}

  \item \textbf{M0: Foundation}
    \begin{itemize}
      \item ระบุกลุ่มเป้าหมาย ขอบเขต และ MVP
      \item กำหนด data model (workspace, project, story, variant, session, task, feedback, round, report)
      \item เลือก architecture: auth, LLM provider, prototype rendering, hosting
      \item สำรวจและเลือก design system และ UI kit
      \item กำหนด requirements ด้าน PDPA และข้อความขอ consent
      \item ตั้งค่า CI/CD, environments และ event analytics
      % \item การกำหนดดัชนีชี้วัดผล และแผนการดำเนินงาน
    \end{itemize}

  \item \textbf{M1: Secure Access}
    \begin{itemize}
      \item จัดทำ Landing page
      \item จัดทำ Register page
      \item การยืนยันตัวตนผ่านอีเมล
      \item การเข้าสู่ระบบพร้อมระบบจัดการข้อผิดพลาด
      \item ระบบรีเซ็ตรหัสผ่าน
      \item จัดทำระบบ routing ตามบทบาทของผู้ใช้
      \item ระบบควบคุมการเข้าถึง
    \end{itemize}

  \item \textbf{M2: Project Setup}
    \begin{itemize}
      \item ระบบสร้าง workspace
      \item ระบบสร้างโปรเจกต์
      \item ฟอร์มกำหนดความต้องการทางการออกแบบ
      \item ระบบจัดเก็บความต้องการทางการออกแบบเพื่อใช้เป็น reusable AI context
    \end{itemize}

  \item \textbf{M3: AI Design Generation}
    \begin{itemize}
      \item การสร้าง User story และ Problem statement
      \item เครื่องมือสร้างคำสั่งที่รวมความต้องการทางการออกแบบและ user story
      \item การสร้าง two-alternative UI
      \item การปรับปรุง UI (Chat-based)
      \item การดูตัวอย่างรูปแบบหน้าจอ
      \item การอนุมัติและมอบหมายงานให้ผู้เข้าร่วมทดสอบ
      \item การเริ่มเซสชันการทดสอบ
      % \item ชุดข้อมูลสำหรับการประเมินคุณภาพผลลัพธ์ของ AI
    \end{itemize}

  \item \textbf{M4: Testing Space}
    \begin{itemize}
      \item การเข้าสู่เซสชันด้วยรหัสหรือลิงก์ พร้อมระบบตรวจสอบความถูกต้องของอายุการใช้งาน
      \item สถานะสิ้นสุดของระบบ
      \item การขอความยินยอมตามกฎหมายคุ้มครองข้อมูลส่วนบุคคล (PDPA) พร้อมระบบบันทึกประวัติการให้ความยินยอม
      \item การกำหนดผู้ใช้งานให้เข้าร่วมการทดสอบ
      \item การแนะนำการใช้งานเครื่องมือทดสอบสำหรับผู้ใช้ใหม่
      \item ระบบทำภารกิจพร้อมการบันทึกพฤติกรรมการโต้ตอบของผู้ใช้
      \item การแสดงความคิดเห็น การระบุตำแหน่งคอมโพเนนต์ การแสดงความคิดเห็นทั้งหน้าจอ หรือการข้ามการแสดงความคิดเห็น
      \item วงจรการทำภารกิจที่เหลืออยู่จนกระทั่งสิ้นสุดการทดสอบของผู้ทดสอบ
      \item การทดสอบข้ามเบราว์เซอร์
    \end{itemize}

  \item \textbf{M5: Feedback}
    \begin{itemize}
      \item การติดตามเซสชันการทดสอบ
      \item การรวบรวมและประมวลผลฟีดแบ็ก
      \item การวิเคราะห์ด้วย AI การหาประเด็นหลัก การติดป้ายกำกับ และการสรุปผล
      \item หน้าจอเพื่อการตรวจสอบของนักออกแบบ การเปลี่ยนป้ายกำกับ การรวม การแยก การจดบันทึก และการทำเครื่องหมาย "ไม่แก้ไข" พร้อมระบุเหตุผล
      \item วงจรการอนุมัติหรือการปรับปรุงแก้ไข
      \item การนำฟีดแบ็กที่เลือกไปป้อนเป็นพรอมต์สำหรับรอบถัดไป
      \item การตัดสินใจสิ้นสุดเซสชัน การทำเวอร์ชันของรอบการทำงาน และประวัติการทำงาน
      \item การทดสอบวงจรการทำงานแบบต้นทางถึงปลายทาง
    \end{itemize}

  \item \textbf{M6: Stakeholder Report}
    \begin{itemize}
      \item รายการโปรเจกต์ (แบบอ่านอย่างเดียว สำหรับงานที่อนุมัติแล้วเท่านั้น)
      \item มุมมองผู้ใช้งานที่แตกต่างกัน
      \item การส่งมอบงานให้นักพัฒนา
      \item สถานะสำหรับผู้จัดการโครงการ
      \item เครื่องมือสร้างรายงาน
      \item การจัดทำรายงานเฉพาะสำหรับแต่ละบทบาท
      \item การส่งออกและการแชร์รายงาน
    \end{itemize}

  \item \textbf{M7: Pilot and Launch}
    \begin{itemize}
      \item ทดสอบ E2E และ UAT
      \item ทดสอบ Security test และ PDPA audit
      \item ทดสอบ Load test ด้วยผู้ทดสอบพร้อมกัน
      \item ทดลองใช้งานกับทีมจริง 2–3 ทีมและเก็บข้อเสนอแนะ
      \item จัดทำคู่มือการใช้งาน การเริ่มต้นใช้งาน และกระบวนการสนับสนุน
      \item ติดตามระบบและแดชบอร์ดตัวชี้วัดทางธุรกิจ
      \item Go live
    \end{itemize}
\end{enumerate}

\subsection{Gantt Chart}
% !TeX root = ../report.tex
% Generated by tools/gantt_to_latex.py from the Gantt sheet. Edit the sheet and regenerate,
% rather than editing this file by hand.
\begingroup
\footnotesize
\setlength{\tabcolsep}{2pt}
\renewcommand{\arraystretch}{1.15}
\begin{longtable}{|>{\RaggedRight\arraybackslash}p{4.4cm}|@{}>{\centering\arraybackslash}p{\dimexpr(\linewidth-4.4cm-1.9cm-4\tabcolsep-30\arrayrulewidth)/27\relax}@{}|@{}>{\centering\arraybackslash}p{\dimexpr(\linewidth-4.4cm-1.9cm-4\tabcolsep-30\arrayrulewidth)/27\relax}@{}|@{}>{\centering\arraybackslash}p{\dimexpr(\linewidth-4.4cm-1.9cm-4\tabcolsep-30\arrayrulewidth)/27\relax}@{}|@{}>{\centering\arraybackslash}p{\dimexpr(\linewidth-4.4cm-1.9cm-4\tabcolsep-30\arrayrulewidth)/27\relax}@{}|@{}>{\centering\arraybackslash}p{\dimexpr(\linewidth-4.4cm-1.9cm-4\tabcolsep-30\arrayrulewidth)/27\relax}@{}|@{}>{\centering\arraybackslash}p{\dimexpr(\linewidth-4.4cm-1.9cm-4\tabcolsep-30\arrayrulewidth)/27\relax}@{}|@{}>{\centering\arraybackslash}p{\dimexpr(\linewidth-4.4cm-1.9cm-4\tabcolsep-30\arrayrulewidth)/27\relax}@{}|@{}>{\centering\arraybackslash}p{\dimexpr(\linewidth-4.4cm-1.9cm-4\tabcolsep-30\arrayrulewidth)/27\relax}@{}|@{}>{\centering\arraybackslash}p{\dimexpr(\linewidth-4.4cm-1.9cm-4\tabcolsep-30\arrayrulewidth)/27\relax}@{}|@{}>{\centering\arraybackslash}p{\dimexpr(\linewidth-4.4cm-1.9cm-4\tabcolsep-30\arrayrulewidth)/27\relax}@{}|@{}>{\centering\arraybackslash}p{\dimexpr(\linewidth-4.4cm-1.9cm-4\tabcolsep-30\arrayrulewidth)/27\relax}@{}|@{}>{\centering\arraybackslash}p{\dimexpr(\linewidth-4.4cm-1.9cm-4\tabcolsep-30\arrayrulewidth)/27\relax}@{}|@{}>{\centering\arraybackslash}p{\dimexpr(\linewidth-4.4cm-1.9cm-4\tabcolsep-30\arrayrulewidth)/27\relax}@{}|@{}>{\centering\arraybackslash}p{\dimexpr(\linewidth-4.4cm-1.9cm-4\tabcolsep-30\arrayrulewidth)/27\relax}@{}|@{}>{\centering\arraybackslash}p{\dimexpr(\linewidth-4.4cm-1.9cm-4\tabcolsep-30\arrayrulewidth)/27\relax}@{}|@{}>{\centering\arraybackslash}p{\dimexpr(\linewidth-4.4cm-1.9cm-4\tabcolsep-30\arrayrulewidth)/27\relax}@{}|@{}>{\centering\arraybackslash}p{\dimexpr(\linewidth-4.4cm-1.9cm-4\tabcolsep-30\arrayrulewidth)/27\relax}@{}|@{}>{\centering\arraybackslash}p{\dimexpr(\linewidth-4.4cm-1.9cm-4\tabcolsep-30\arrayrulewidth)/27\relax}@{}|@{}>{\centering\arraybackslash}p{\dimexpr(\linewidth-4.4cm-1.9cm-4\tabcolsep-30\arrayrulewidth)/27\relax}@{}|@{}>{\centering\arraybackslash}p{\dimexpr(\linewidth-4.4cm-1.9cm-4\tabcolsep-30\arrayrulewidth)/27\relax}@{}|@{}>{\centering\arraybackslash}p{\dimexpr(\linewidth-4.4cm-1.9cm-4\tabcolsep-30\arrayrulewidth)/27\relax}@{}|@{}>{\centering\arraybackslash}p{\dimexpr(\linewidth-4.4cm-1.9cm-4\tabcolsep-30\arrayrulewidth)/27\relax}@{}|@{}>{\centering\arraybackslash}p{\dimexpr(\linewidth-4.4cm-1.9cm-4\tabcolsep-30\arrayrulewidth)/27\relax}@{}|@{}>{\centering\arraybackslash}p{\dimexpr(\linewidth-4.4cm-1.9cm-4\tabcolsep-30\arrayrulewidth)/27\relax}@{}|@{}>{\centering\arraybackslash}p{\dimexpr(\linewidth-4.4cm-1.9cm-4\tabcolsep-30\arrayrulewidth)/27\relax}@{}|@{}>{\centering\arraybackslash}p{\dimexpr(\linewidth-4.4cm-1.9cm-4\tabcolsep-30\arrayrulewidth)/27\relax}@{}|@{}>{\centering\arraybackslash}p{\dimexpr(\linewidth-4.4cm-1.9cm-4\tabcolsep-30\arrayrulewidth)/27\relax}@{}|>{\RaggedRight\arraybackslash}p{1.9cm}|}
\caption{แผนการดำเนินงานในภาคการศึกษาที่ 1}\label{tab:gantt-semester-1}\\
\hline
\multirow{2}{*}{\textbf{รายละเอียดงาน}} & \multicolumn{5}{c|}{\textbf{ก.ค. 69}} & \multicolumn{4}{c|}{\textbf{ส.ค. 69}} & \multicolumn{5}{c|}{\textbf{ก.ย. 69}} & \multicolumn{4}{c|}{\textbf{ต.ค. 69}} & \multicolumn{4}{c|}{\textbf{พ.ย. 69}} & \multicolumn{5}{c|}{\textbf{ธ.ค. 69}} & \multirow{2}{*}{\textbf{ผู้รับผิดชอบ}} \\
\cline{2-28}
 & 1 & 2 & 3 & 4 & 5 & 1 & 2 & 3 & 4 & 1 & 2 & 3 & 4 & 5 & 1 & 2 & 3 & 4 & 1 & 2 & 3 & 4 & 1 & 2 & 3 & 4 & 5 & \\
\hline
\endfirsthead
\multicolumn{29}{l}{\textit{ตารางที่ \thetable\ (ต่อ)}}\\
\hline
\multirow{2}{*}{\textbf{รายละเอียดงาน}} & \multicolumn{5}{c|}{\textbf{ก.ค. 69}} & \multicolumn{4}{c|}{\textbf{ส.ค. 69}} & \multicolumn{5}{c|}{\textbf{ก.ย. 69}} & \multicolumn{4}{c|}{\textbf{ต.ค. 69}} & \multicolumn{4}{c|}{\textbf{พ.ย. 69}} & \multicolumn{5}{c|}{\textbf{ธ.ค. 69}} & \multirow{2}{*}{\textbf{ผู้รับผิดชอบ}} \\
\cline{2-28}
 & 1 & 2 & 3 & 4 & 5 & 1 & 2 & 3 & 4 & 1 & 2 & 3 & 4 & 5 & 1 & 2 & 3 & 4 & 1 & 2 & 3 & 4 & 1 & 2 & 3 & 4 & 5 & \\
\hline
\endhead
\multicolumn{29}{|l|}{\textbf{Documentation}} \\ \hline
ศึกษาค้นคว้าหัวข้อโครงงาน & \cellcolor{ganttbar} & \cellcolor{ganttbar} & \cellcolor{ganttbar} & \cellcolor{ganttbar} & \cellcolor{ganttbar} & \cellcolor{ganttbar} & \cellcolor{ganttbar} &  &  &  &  &  &  &  &  &  &  &  &  &  &  &  &  &  &  &  &  & ทุกคน \\ \hline
จัดทำเอกสาร Project Idea &  &  &  &  &  & \cellcolor{ganttbar} & \cellcolor{ganttbar} & \cellcolor{ganttbar} & \cellcolor{ganttbar} &  &  &  &  &  &  &  &  &  &  &  &  &  &  &  &  &  &  & ทุกคน \\ \hline
เอกสารและงานวิจัยที่เกี่ยวข้อง &  &  &  &  &  &  &  &  &  & \cellcolor{ganttbar} & \cellcolor{ganttbar} & \cellcolor{ganttbar} &  &  &  &  &  &  &  &  &  &  &  &  &  &  &  & ภูบดินทร์, นาวิน \\ \hline
จัดทำเอกสาร Proposal report &  &  &  &  &  &  &  &  &  & \cellcolor{ganttbar} & \cellcolor{ganttbar} & \cellcolor{ganttbar} & \cellcolor{ganttbar} & \cellcolor{ganttbar} &  &  &  &  &  &  &  &  &  &  &  &  &  & ทุกคน \\ \hline
นำเสนอ Proposal &  &  &  &  &  &  &  &  &  &  &  &  &  & \cellcolor{ganttbar} & \cellcolor{ganttbar} &  &  &  &  &  &  &  &  &  &  &  &  & ทุกคน \\ \hline
จัดทำเอกสาร Progress report ภาคเรียนที่ 1 &  &  &  &  &  &  &  &  &  &  &  &  &  &  &  &  &  &  &  & \cellcolor{ganttbar} & \cellcolor{ganttbar} & \cellcolor{ganttbar} & \cellcolor{ganttbar} &  &  &  &  & ทุกคน \\ \hline
นำเสนอ Progress ภาคเรียนที่ 1 &  &  &  &  &  &  &  &  &  &  &  &  &  &  &  &  &  &  &  &  &  &  & \cellcolor{ganttbar} & \cellcolor{ganttbar} &  &  &  & ทุกคน \\ \hline
\multicolumn{29}{|l|}{\textbf{Requirements Analysis \& System Design}} \\ \hline
สัมภาษณ์ผู้มีส่วนได้ส่วนเสีย &  &  &  &  &  &  &  &  &  & \cellcolor{ganttbar} & \cellcolor{ganttbar} &  &  &  &  &  &  &  &  &  &  &  &  &  &  &  &  & กัญญ์ญาพัชร, ปัณณวิชญ์ \\ \hline
ระบุ Requirements &  &  &  &  &  &  &  &  &  &  & \cellcolor{ganttbar} & \cellcolor{ganttbar} &  &  &  &  &  &  &  &  &  &  &  &  &  &  &  & ปัณณวิชญ์, กัญญ์ญาพัชร \\ \hline
ระบุ Functional \& Non-Functional Requirements &  &  &  &  &  &  &  &  &  &  & \cellcolor{ganttbar} & \cellcolor{ganttbar} &  &  &  &  &  &  &  &  &  &  &  &  &  &  &  & กัญญ์ญาพัชร, ปัณณวิชญ์ \\ \hline
จัดทำ Use Cases \& Sequence Diagram &  &  &  &  &  &  &  &  &  &  &  &  & \cellcolor{ganttbar} &  &  &  &  &  &  &  &  &  &  &  &  &  &  & นาวิน, กัญญ์ญาพัชร \\ \hline
วางโครงสร้าง System Architecture &  &  &  &  &  &  &  &  &  &  &  & \cellcolor{ganttbar} & \cellcolor{ganttbar} &  &  &  &  &  &  &  &  &  &  &  &  &  &  & ภูบดินทร์ \\ \hline
วางโครงสร้าง AI Workflow \& Prompt Architecture &  &  &  &  &  &  &  &  &  &  &  &  &  &  &  & \cellcolor{ganttbar} & \cellcolor{ganttbar} &  &  &  &  &  &  &  &  &  &  & ภูบดินทร์, ปัณณวิชญ์ \\ \hline
ออกแบบ Database &  &  &  &  &  &  &  &  &  &  &  &  &  &  &  & \cellcolor{ganttbar} & \cellcolor{ganttbar} &  &  &  &  &  &  &  &  &  &  & ภูบดินทร์, นาวิน \\ \hline
ออกแบบ API Design &  &  &  &  &  &  &  &  &  &  &  &  &  &  &  &  & \cellcolor{ganttbar} & \cellcolor{ganttbar} &  &  &  &  &  &  &  &  &  & ภูบดินทร์ \\ \hline
ออกแบบ UI/UX Design (User Flow \& Wireframe) &  &  &  &  &  &  &  &  &  &  &  & \cellcolor{ganttbar} &  &  &  &  &  &  &  &  &  &  &  &  &  &  &  & ปัณณวิชญ์ \\ \hline
System prototype (PoC) &  &  &  &  &  &  &  &  &  &  &  &  &  &  &  &  & \cellcolor{ganttbar} & \cellcolor{ganttbar} &  &  &  &  &  &  &  &  &  & ทุกคน \\ \hline
\multicolumn{29}{|l|}{\textbf{M0: Foundation}} \\ \hline
ระบุกลุ่มเป้าหมาย ขอบเขต และ MVP &  &  &  &  &  &  &  &  &  &  & \cellcolor{ganttbar} & \cellcolor{ganttbar} &  &  &  &  &  &  &  &  &  &  &  &  &  &  &  & กัญญ์ญาพัชร, ปัณณวิชญ์ \\ \hline
กำหนด data model (workspace, project, story, variant, session, task, feedback, round, report) &  &  &  &  &  &  &  &  &  &  &  &  &  &  &  & \cellcolor{ganttbar} & \cellcolor{ganttbar} &  &  &  &  &  &  &  &  &  &  & นาวิน, ภูบดินทร์ \\ \hline
เลือก architecture: auth, LLM provider, prototype rendering, hosting &  &  &  &  &  &  &  &  &  &  &  &  &  &  &  & \cellcolor{ganttbar} & \cellcolor{ganttbar} &  &  &  &  &  &  &  &  &  &  & ภูบดินทร์, นาวิน \\ \hline
สำรวจและเลือก design system และ UI kit &  &  &  &  &  &  &  &  &  &  &  &  &  &  &  &  & \cellcolor{ganttbar} & \cellcolor{ganttbar} &  &  &  &  &  &  &  &  &  & ปัณณวิชญ์ \\ \hline
กำหนด requirements ด้าน PDPA และข้อความขอ consent &  &  &  &  &  &  &  &  &  &  &  &  &  &  &  & \cellcolor{ganttbar} & \cellcolor{ganttbar} &  &  &  &  &  &  &  &  &  &  & กัญญ์ญาพัชร \\ \hline
ตั้งค่า CI/CD, environments และ event analytics &  &  &  &  &  &  &  &  &  &  &  &  &  &  &  &  & \cellcolor{ganttbar} & \cellcolor{ganttbar} &  &  &  &  &  &  &  &  &  & ภูบดินทร์ \\ \hline
การกำหนดดัชนีชี้วัดผล และแผนการดำเนินงาน &  &  &  &  &  &  &  &  &  &  &  &  &  &  &  & \cellcolor{ganttbar} &  &  &  &  &  &  &  &  &  &  &  & ทุกคน \\ \hline
\multicolumn{29}{|l|}{\textbf{M1: Secure Access}} \\ \hline
จัดทำ Landing page &  &  &  &  &  &  &  &  &  &  &  &  &  &  &  &  & \cellcolor{ganttbar} & \cellcolor{ganttbar} &  &  &  &  &  &  &  &  &  & -- \\ \hline
จัดทำ Register page &  &  &  &  &  &  &  &  &  &  &  &  &  &  &  &  & \cellcolor{ganttbar} & \cellcolor{ganttbar} & \cellcolor{ganttbar} &  &  &  &  &  &  &  &  & -- \\ \hline
การยืนยันตัวตนผ่านอีเมล &  &  &  &  &  &  &  &  &  &  &  &  &  &  &  &  &  &  & \cellcolor{ganttbar} & \cellcolor{ganttbar} &  &  &  &  &  &  &  & -- \\ \hline
การเข้าสู่ระบบพร้อมระบบจัดการข้อผิดพลาด &  &  &  &  &  &  &  &  &  &  &  &  &  &  &  &  &  & \cellcolor{ganttbar} & \cellcolor{ganttbar} & \cellcolor{ganttbar} &  &  &  &  &  &  &  & -- \\ \hline
ระบบรีเซ็ตรหัสผ่าน &  &  &  &  &  &  &  &  &  &  &  &  &  &  &  &  &  &  &  & \cellcolor{ganttbar} & \cellcolor{ganttbar} & \cellcolor{ganttbar} &  &  &  &  &  & -- \\ \hline
ระบบควบคุมการเข้าถึง &  &  &  &  &  &  &  &  &  &  &  &  &  &  &  &  &  &  &  & \cellcolor{ganttbar} & \cellcolor{ganttbar} & \cellcolor{ganttbar} &  &  &  &  &  & -- \\ \hline
จัดทำระบบ routing ตามบทบาทของผู้ใช้ &  &  &  &  &  &  &  &  &  &  &  &  &  &  &  &  &  &  &  &  &  & \cellcolor{ganttbar} & \cellcolor{ganttbar} & \cellcolor{ganttbar} &  &  &  & -- \\ \hline
ทดสอบระบบยืนยันตัวตนและการควบคุมการเข้าถึง &  &  &  &  &  &  &  &  &  &  &  &  &  &  &  &  &  &  &  &  &  & \cellcolor{ganttbar} & \cellcolor{ganttbar} & \cellcolor{ganttbar} &  &  &  & -- \\ \hline
\multicolumn{29}{|l|}{\textbf{M2: Project Setup}} \\ \hline
ระบบสร้าง workspace &  &  &  &  &  &  &  &  &  &  &  &  &  &  &  &  &  &  &  &  &  & \cellcolor{ganttbar} & \cellcolor{ganttbar} & \cellcolor{ganttbar} &  &  &  & -- \\ \hline
ระบบสร้างโปรเจกต์ &  &  &  &  &  &  &  &  &  &  &  &  &  &  &  &  &  &  &  &  &  &  & \cellcolor{ganttbar} & \cellcolor{ganttbar} &  &  &  & -- \\ \hline
ฟอร์มกำหนดความต้องการทางการออกแบบ &  &  &  &  &  &  &  &  &  &  &  &  &  &  &  &  &  &  &  &  &  &  & \cellcolor{ganttbar} & \cellcolor{ganttbar} & \cellcolor{ganttbar} & \cellcolor{ganttbar} &  & -- \\ \hline
ระบบจัดเก็บความต้องการทางการออกแบบเพื่อใช้เป็น reusable AI context &  &  &  &  &  &  &  &  &  &  &  &  &  &  &  &  &  &  &  &  &  &  &  & \cellcolor{ganttbar} & \cellcolor{ganttbar} & \cellcolor{ganttbar} &  & -- \\ \hline
\multicolumn{29}{|l|}{\textbf{M3: AI Design Generation}} \\ \hline
การสร้าง User story และ Problem statement &  &  &  &  &  &  &  &  &  &  &  &  &  &  &  &  &  &  &  &  &  &  &  & \cellcolor{ganttbar} & \cellcolor{ganttbar} & \cellcolor{ganttbar} &  & -- \\ \hline
เครื่องมือสร้างคำสั่งที่รวมความต้องการทางการออกแบบและ user story &  &  &  &  &  &  &  &  &  &  &  &  &  &  &  &  &  &  &  &  &  &  &  &  & \cellcolor{ganttbar} & \cellcolor{ganttbar} & \cellcolor{ganttbar} & -- \\ \hline
การสร้าง two-alternative UI &  &  &  &  &  &  &  &  &  &  &  &  &  &  &  &  &  &  &  &  &  &  &  &  &  & \cellcolor{ganttbar} & \cellcolor{ganttbar} & -- \\ \hline
การปรับปรุง UI (Chat-based) &  &  &  &  &  &  &  &  &  &  &  &  &  &  &  &  &  &  &  &  &  &  &  &  &  &  & \cellcolor{ganttbar} & -- \\ \hline
การดูตัวอย่างรูปแบบหน้าจอ &  &  &  &  &  &  &  &  &  &  &  &  &  &  &  &  &  &  &  &  &  &  &  &  &  & \cellcolor{ganttbar} & \cellcolor{ganttbar} & -- \\ \hline
ชุดข้อมูลสำหรับการประเมินคุณภาพผลลัพธ์ของ AI &  &  &  &  &  &  &  &  &  &  &  &  &  &  &  &  &  &  &  &  &  &  &  &  &  &  & \cellcolor{ganttbar} & -- \\ \hline
\end{longtable}
\endgroup

\begingroup
\footnotesize
\setlength{\tabcolsep}{2pt}
\renewcommand{\arraystretch}{1.15}
\begin{longtable}{|>{\RaggedRight\arraybackslash}p{4.4cm}|@{}>{\centering\arraybackslash}p{\dimexpr(\linewidth-4.4cm-1.9cm-4\tabcolsep-24\arrayrulewidth)/21\relax}@{}|@{}>{\centering\arraybackslash}p{\dimexpr(\linewidth-4.4cm-1.9cm-4\tabcolsep-24\arrayrulewidth)/21\relax}@{}|@{}>{\centering\arraybackslash}p{\dimexpr(\linewidth-4.4cm-1.9cm-4\tabcolsep-24\arrayrulewidth)/21\relax}@{}|@{}>{\centering\arraybackslash}p{\dimexpr(\linewidth-4.4cm-1.9cm-4\tabcolsep-24\arrayrulewidth)/21\relax}@{}|@{}>{\centering\arraybackslash}p{\dimexpr(\linewidth-4.4cm-1.9cm-4\tabcolsep-24\arrayrulewidth)/21\relax}@{}|@{}>{\centering\arraybackslash}p{\dimexpr(\linewidth-4.4cm-1.9cm-4\tabcolsep-24\arrayrulewidth)/21\relax}@{}|@{}>{\centering\arraybackslash}p{\dimexpr(\linewidth-4.4cm-1.9cm-4\tabcolsep-24\arrayrulewidth)/21\relax}@{}|@{}>{\centering\arraybackslash}p{\dimexpr(\linewidth-4.4cm-1.9cm-4\tabcolsep-24\arrayrulewidth)/21\relax}@{}|@{}>{\centering\arraybackslash}p{\dimexpr(\linewidth-4.4cm-1.9cm-4\tabcolsep-24\arrayrulewidth)/21\relax}@{}|@{}>{\centering\arraybackslash}p{\dimexpr(\linewidth-4.4cm-1.9cm-4\tabcolsep-24\arrayrulewidth)/21\relax}@{}|@{}>{\centering\arraybackslash}p{\dimexpr(\linewidth-4.4cm-1.9cm-4\tabcolsep-24\arrayrulewidth)/21\relax}@{}|@{}>{\centering\arraybackslash}p{\dimexpr(\linewidth-4.4cm-1.9cm-4\tabcolsep-24\arrayrulewidth)/21\relax}@{}|@{}>{\centering\arraybackslash}p{\dimexpr(\linewidth-4.4cm-1.9cm-4\tabcolsep-24\arrayrulewidth)/21\relax}@{}|@{}>{\centering\arraybackslash}p{\dimexpr(\linewidth-4.4cm-1.9cm-4\tabcolsep-24\arrayrulewidth)/21\relax}@{}|@{}>{\centering\arraybackslash}p{\dimexpr(\linewidth-4.4cm-1.9cm-4\tabcolsep-24\arrayrulewidth)/21\relax}@{}|@{}>{\centering\arraybackslash}p{\dimexpr(\linewidth-4.4cm-1.9cm-4\tabcolsep-24\arrayrulewidth)/21\relax}@{}|@{}>{\centering\arraybackslash}p{\dimexpr(\linewidth-4.4cm-1.9cm-4\tabcolsep-24\arrayrulewidth)/21\relax}@{}|@{}>{\centering\arraybackslash}p{\dimexpr(\linewidth-4.4cm-1.9cm-4\tabcolsep-24\arrayrulewidth)/21\relax}@{}|@{}>{\centering\arraybackslash}p{\dimexpr(\linewidth-4.4cm-1.9cm-4\tabcolsep-24\arrayrulewidth)/21\relax}@{}|@{}>{\centering\arraybackslash}p{\dimexpr(\linewidth-4.4cm-1.9cm-4\tabcolsep-24\arrayrulewidth)/21\relax}@{}|@{}>{\centering\arraybackslash}p{\dimexpr(\linewidth-4.4cm-1.9cm-4\tabcolsep-24\arrayrulewidth)/21\relax}@{}|>{\RaggedRight\arraybackslash}p{1.9cm}|}
\caption{แผนการดำเนินงานในภาคการศึกษาที่ 2}\label{tab:gantt-semester-2}\\
\hline
\multirow{2}{*}{\textbf{รายละเอียดงาน}} & \multicolumn{4}{c|}{\textbf{ม.ค. 70}} & \multicolumn{4}{c|}{\textbf{ก.พ. 70}} & \multicolumn{5}{c|}{\textbf{มี.ค. 70}} & \multicolumn{4}{c|}{\textbf{เม.ย. 70}} & \multicolumn{4}{c|}{\textbf{พ.ค. 70}} & \multirow{2}{*}{\textbf{ผู้รับผิดชอบ}} \\
\cline{2-22}
 & 1 & 2 & 3 & 4 & 1 & 2 & 3 & 4 & 1 & 2 & 3 & 4 & 5 & 1 & 2 & 3 & 4 & 1 & 2 & 3 & 4 & \\
\hline
\endfirsthead
\multicolumn{23}{l}{\textit{ตารางที่ \thetable\ (ต่อ)}}\\
\hline
\multirow{2}{*}{\textbf{รายละเอียดงาน}} & \multicolumn{4}{c|}{\textbf{ม.ค. 70}} & \multicolumn{4}{c|}{\textbf{ก.พ. 70}} & \multicolumn{5}{c|}{\textbf{มี.ค. 70}} & \multicolumn{4}{c|}{\textbf{เม.ย. 70}} & \multicolumn{4}{c|}{\textbf{พ.ค. 70}} & \multirow{2}{*}{\textbf{ผู้รับผิดชอบ}} \\
\cline{2-22}
 & 1 & 2 & 3 & 4 & 1 & 2 & 3 & 4 & 1 & 2 & 3 & 4 & 5 & 1 & 2 & 3 & 4 & 1 & 2 & 3 & 4 & \\
\hline
\endhead
\multicolumn{23}{|l|}{\textbf{Documentation}} \\ \hline
จัดทำเอกสาร Final Report &  &  &  &  &  &  &  &  &  &  &  &  & \cellcolor{ganttbar} & \cellcolor{ganttbar} & \cellcolor{ganttbar} & \cellcolor{ganttbar} & \cellcolor{ganttbar} &  &  &  &  & ทุกคน \\ \hline
นำเสนอ Final Defense &  &  &  &  &  &  &  &  &  &  &  &  & \cellcolor{ganttbar} & \cellcolor{ganttbar} & \cellcolor{ganttbar} & \cellcolor{ganttbar} & \cellcolor{ganttbar} &  &  &  &  & ทุกคน \\ \hline
\multicolumn{23}{|l|}{\textbf{M3: AI Design Generation}} \\ \hline
เครื่องมือสร้างคำสั่งที่รวมความต้องการทางการออกแบบและ user story & \cellcolor{ganttbar} &  &  &  &  &  &  &  &  &  &  &  &  &  &  &  &  &  &  &  &  & -- \\ \hline
การสร้าง two-alternative UI & \cellcolor{ganttbar} & \cellcolor{ganttbar} &  &  &  &  &  &  &  &  &  &  &  &  &  &  &  &  &  &  &  & -- \\ \hline
การปรับปรุง UI (Chat-based) & \cellcolor{ganttbar} & \cellcolor{ganttbar} &  &  &  &  &  &  &  &  &  &  &  &  &  &  &  &  &  &  &  & -- \\ \hline
การดูตัวอย่างรูปแบบหน้าจอ & \cellcolor{ganttbar} & \cellcolor{ganttbar} &  &  &  &  &  &  &  &  &  &  &  &  &  &  &  &  &  &  &  & -- \\ \hline
การอนุมัติและมอบหมายงานให้ผู้เข้าร่วมทดสอบ &  & \cellcolor{ganttbar} & \cellcolor{ganttbar} &  &  &  &  &  &  &  &  &  &  &  &  &  &  &  &  &  &  & -- \\ \hline
การเริ่มเซสชันการทดสอบ &  & \cellcolor{ganttbar} & \cellcolor{ganttbar} &  &  &  &  &  &  &  &  &  &  &  &  &  &  &  &  &  &  & -- \\ \hline
ชุดข้อมูลสำหรับการประเมินคุณภาพผลลัพธ์ของ AI & \cellcolor{ganttbar} & \cellcolor{ganttbar} & \cellcolor{ganttbar} &  &  &  &  &  &  &  &  &  &  &  &  &  &  &  &  &  &  & -- \\ \hline
\multicolumn{23}{|l|}{\textbf{M4: Testing Space}} \\ \hline
การเข้าสู่เซสชันด้วยรหัสหรือลิงก์ พร้อมระบบตรวจสอบความถูกต้องของอายุการใช้งาน &  &  &  &  & \cellcolor{ganttbar} & \cellcolor{ganttbar} &  &  &  &  &  &  &  &  &  &  &  &  &  &  &  & -- \\ \hline
สถานะสิ้นสุดของระบบ &  &  &  &  &  & \cellcolor{ganttbar} & \cellcolor{ganttbar} &  &  &  &  &  &  &  &  &  &  &  &  &  &  & -- \\ \hline
การขอความยินยอมตามกฎหมายคุ้มครองข้อมูลส่วนบุคคล (PDPA) พร้อมระบบบันทึกประวัติการให้ความยินยอม &  &  &  &  &  &  & \cellcolor{ganttbar} & \cellcolor{ganttbar} &  &  &  &  &  &  &  &  &  &  &  &  &  & -- \\ \hline
การกำหนดผู้ใช้งานให้เข้าร่วมการทดสอบ &  &  &  &  &  &  &  & \cellcolor{ganttbar} & \cellcolor{ganttbar} &  &  &  &  &  &  &  &  &  &  &  &  & -- \\ \hline
การแนะนำการใช้งานเครื่องมือทดสอบสำหรับผู้ใช้ใหม่ &  &  &  &  &  &  & \cellcolor{ganttbar} & \cellcolor{ganttbar} &  &  &  &  &  &  &  &  &  &  &  &  &  & -- \\ \hline
ระบบทำภารกิจพร้อมการบันทึกพฤติกรรมการโต้ตอบของผู้ใช้ &  &  &  &  &  &  & \cellcolor{ganttbar} & \cellcolor{ganttbar} &  &  &  &  &  &  &  &  &  &  &  &  &  & -- \\ \hline
การทบทวนหน้าจอหลังทำภารกิจและคำถามเฉพาะจุด &  &  &  &  &  &  &  & \cellcolor{ganttbar} & \cellcolor{ganttbar} & \cellcolor{ganttbar} &  &  &  &  &  &  &  &  &  &  &  & -- \\ \hline
การแสดงความคิดเห็น การระบุตำแหน่งคอมโพเนนต์ การแสดงความคิดเห็นทั้งหน้าจอ หรือการข้ามการแสดงความคิดเห็น &  &  &  &  &  &  &  & \cellcolor{ganttbar} & \cellcolor{ganttbar} & \cellcolor{ganttbar} & \cellcolor{ganttbar} &  &  &  &  &  &  &  &  &  &  & -- \\ \hline
วงจรการทำภารกิจที่เหลืออยู่จนกระทั่งสิ้นสุดการทดสอบของผู้ทดสอบ &  &  &  &  &  &  &  &  & \cellcolor{ganttbar} & \cellcolor{ganttbar} &  &  &  &  &  &  &  &  &  &  &  & -- \\ \hline
การทดสอบข้ามเบราว์เซอร์ &  &  &  &  &  &  &  &  &  & \cellcolor{ganttbar} & \cellcolor{ganttbar} & \cellcolor{ganttbar} &  &  &  &  &  &  &  &  &  & -- \\ \hline
\multicolumn{23}{|l|}{\textbf{M5: Feedback}} \\ \hline
การติดตามเซสชันการทดสอบ &  &  &  &  &  &  &  &  &  & \cellcolor{ganttbar} & \cellcolor{ganttbar} & \cellcolor{ganttbar} &  &  &  &  &  &  &  &  &  & -- \\ \hline
การรวบรวมและประมวลผลฟีดแบ็ก &  &  &  &  &  &  &  &  &  & \cellcolor{ganttbar} & \cellcolor{ganttbar} & \cellcolor{ganttbar} &  &  &  &  &  &  &  &  &  & -- \\ \hline
การวิเคราะห์ด้วย AI การหาประเด็นหลัก การติดป้ายกำกับ และการสรุปผล &  &  &  &  &  &  &  &  &  &  & \cellcolor{ganttbar} & \cellcolor{ganttbar} & \cellcolor{ganttbar} & \cellcolor{ganttbar} &  &  &  &  &  &  &  & -- \\ \hline
หน้าจอเพื่อการตรวจสอบของนักออกแบบ การเปลี่ยนป้ายกำกับ การรวม การแยก การจดบันทึก และการทำเครื่องหมาย "ไม่แก้ไข" พร้อมระบุเหตุผล &  &  &  &  &  &  &  &  &  &  & \cellcolor{ganttbar} & \cellcolor{ganttbar} & \cellcolor{ganttbar} & \cellcolor{ganttbar} &  &  &  &  &  &  &  & -- \\ \hline
วงจรการอนุมัติหรือการปรับปรุงแก้ไข &  &  &  &  &  &  &  &  &  &  &  & \cellcolor{ganttbar} & \cellcolor{ganttbar} &  &  &  &  &  &  &  &  & -- \\ \hline
การนำฟีดแบ็กที่เลือกไปป้อนเป็นพรอมต์สำหรับรอบถัดไป &  &  &  &  &  &  &  &  &  &  &  & \cellcolor{ganttbar} & \cellcolor{ganttbar} & \cellcolor{ganttbar} &  &  &  &  &  &  &  & -- \\ \hline
การตัดสินใจสิ้นสุดเซสชัน การทำเวอร์ชันของรอบการทำงาน และประวัติการทำงาน &  &  &  &  &  &  &  &  &  & \cellcolor{ganttbar} & \cellcolor{ganttbar} & \cellcolor{ganttbar} &  &  &  &  &  &  &  &  &  & -- \\ \hline
การทดสอบวงจรการทำงานแบบต้นทางถึงปลายทาง &  &  &  &  &  &  &  &  &  &  &  & \cellcolor{ganttbar} & \cellcolor{ganttbar} &  &  &  &  &  &  &  &  & -- \\ \hline
\multicolumn{23}{|l|}{\textbf{M6: Stakeholder Report}} \\ \hline
รายการโปรเจกต์ (แบบอ่านอย่างเดียว สำหรับงานที่อนุมัติแล้วเท่านั้น) &  &  &  &  &  &  &  &  &  &  &  & \cellcolor{ganttbar} & \cellcolor{ganttbar} &  &  &  &  &  &  &  &  & -- \\ \hline
มุมมองผู้ใช้งานที่แตกต่างกัน &  &  &  &  &  &  &  &  &  &  &  &  &  & \cellcolor{ganttbar} & \cellcolor{ganttbar} & \cellcolor{ganttbar} &  &  &  &  &  & -- \\ \hline
การส่งมอบงานให้นักพัฒนา &  &  &  &  &  &  &  &  &  &  &  &  & \cellcolor{ganttbar} & \cellcolor{ganttbar} & \cellcolor{ganttbar} & \cellcolor{ganttbar} &  &  &  &  &  & -- \\ \hline
สถานะสำหรับผู้จัดการโครงการ &  &  &  &  &  &  &  &  &  &  &  &  & \cellcolor{ganttbar} & \cellcolor{ganttbar} & \cellcolor{ganttbar} &  &  &  &  &  &  & -- \\ \hline
เครื่องมือสร้างรายงาน &  &  &  &  &  &  &  &  &  &  &  &  & \cellcolor{ganttbar} & \cellcolor{ganttbar} & \cellcolor{ganttbar} & \cellcolor{ganttbar} &  &  &  &  &  & -- \\ \hline
การจัดทำรายงานเฉพาะสำหรับแต่ละบทบาท &  &  &  &  &  &  &  &  &  &  &  &  &  & \cellcolor{ganttbar} & \cellcolor{ganttbar} & \cellcolor{ganttbar} & \cellcolor{ganttbar} &  &  &  &  & -- \\ \hline
การส่งออกและการแชร์รายงาน &  &  &  &  &  &  &  &  &  &  &  &  &  &  &  & \cellcolor{ganttbar} & \cellcolor{ganttbar} &  &  &  &  & -- \\ \hline
\multicolumn{23}{|l|}{\textbf{M7: Pilot and Launch}} \\ \hline
ทดสอบ E2E และ UAT &  &  &  &  &  &  &  &  &  &  &  &  &  &  &  &  & \cellcolor{ganttbar} & \cellcolor{ganttbar} &  &  &  & -- \\ \hline
ทดสอบ Security test และ PDPA audit &  &  &  &  &  &  &  &  &  &  &  &  &  &  &  &  &  & \cellcolor{ganttbar} &  &  &  & -- \\ \hline
ทดสอบ Load test ด้วยผู้ทดสอบพร้อมกัน &  &  &  &  &  &  &  &  &  &  &  &  &  &  &  &  &  & \cellcolor{ganttbar} &  &  &  & -- \\ \hline
ทดลองใช้งานกับทีมจริง 2–3 ทีมและเก็บข้อเสนอแนะ &  &  &  &  &  &  &  &  &  &  &  &  &  &  &  &  &  & \cellcolor{ganttbar} & \cellcolor{ganttbar} & \cellcolor{ganttbar} & \cellcolor{ganttbar} & -- \\ \hline
จัดทำคู่มือการใช้งาน การเริ่มต้นใช้งาน และกระบวนการสนับสนุน &  &  &  &  &  &  &  &  &  &  &  &  &  &  &  &  & \cellcolor{ganttbar} & \cellcolor{ganttbar} &  &  &  & -- \\ \hline
ติดตามระบบและแดชบอร์ดตัวชี้วัดทางธุรกิจ &  &  &  &  &  &  &  &  &  &  &  &  &  &  &  &  & \cellcolor{ganttbar} & \cellcolor{ganttbar} &  &  &  & -- \\ \hline
Go live &  &  &  &  &  &  &  &  &  &  &  &  &  &  &  &  &  &  &  & \cellcolor{ganttbar} & \cellcolor{ganttbar} & -- \\ \hline
\end{longtable}
\endgroup



\section{ผลการดำเนินงานที่ส่งมอบ}
\subsection{ผลการดำเนินงานในภาคการศึกษาที่ 1}
\begin{enumerate}
  \item \textbf{M0: Foundation}
    \begin{itemize}
      \item Software Requirement Specification, Data Model, System Architecture, Design System \& UI Kit, PDPA \& Consent, CI/CD, Sprint planning
    \end{itemize}
  \item \textbf{M1: Secure Access}
    \begin{itemize}
      \item ระบบสามารถสมัครสมาชิกพร้อมยืนยันอีเมล, เข้าสู่ระบบ ลืมและรีเซ็ตรหัสผ่าน ระบบสิทธิ์ตามบทบาท (Designer, Client, Developer, PM)
    \end{itemize}
  \item \textbf{M2: Project Setup}
    \begin{itemize}
      \item ระบบสามารถสร้าง workspace และโปรเจกต์, แบบฟอร์มความต้องการทางการออกแบบ (description, persona, design guideline, CI, tone of voice)
    \end{itemize}
  \item \textbf{M3: AI Design Generation}
    \begin{itemize}
      \item ระบบสามารถสร้าง user story และ problem statement สร้าง two-alternative UI ด้วย Generative AI, Chat-based UI refinement, prototype preview อนุมัติและมอบหมายงานทดสอบ สร้าง test session (โดยสามารถเข้าถึงด้วยรหัสหรือลิงก์)
    \end{itemize}
\end{enumerate}

\subsection{ผลการดำเนินงานในภาคการศึกษาที่ 2}
\begin{enumerate}
  \item \textbf{M4: Testing Space}
    \begin{itemize}
      \item ผู้เข้าร่วมทดสอบสามารถเข้า session ด้วยรหัสหรือลิงก์ ทำ task พร้อมให้ข้อเสนอแนะเชิงบริบท
    \end{itemize}
  \item \textbf{M5: Feedback}
    \begin{itemize}
      \item ระบบสามารถรวบรวมข้อเสนอแนะ AI สามารถวิเคราะห์รูปแบบ ข้อมูลเชิงลึกและประเด็นสำคัญ Designer ตรวจทานและแก้ไขผลการวิเคราะห์ เลือกข้อเสนอแนะที่จะปรับปรุงในรอบถัดไป บันทึกประวัติแต่ละรอบ
    \end{itemize}
  \item \textbf{M6: Stakeholder Report}
    \begin{itemize}
      \item ระบบสามารถแสดงหน้ารายการโปรเจกต์ (อ่านอย่างเดียว) มุมมองของ Client, Developer และ PM, รายงานสำหรับแต่ละบทบาท พร้อมส่งออก PDF หรือลิงก์
    \end{itemize}
  \item \textbf{M7: Pilot and Launch}
    \begin{itemize}
      \item ทดสอบระบบ (UAT, Load test, Security Test) ทดลองใช้กับ 2–3 ทีมตัวอย่าง, เอกสารการประกอบการใช้งาน เปิดใช้งานจริง
    \end{itemize}
\end{enumerate}

